\chapter{光谱分类案例}

\section{案例目标}

这一章把前面关于数据预处理、降维、分类和模型诊断的内容重新装配成一个完整的小型科学项目。目标是让读者真正体验到：面对一批一维光谱，怎样一步一步走到一个可解释的分类结果。

\begin{enumerate}
    \item 理解为什么光谱分类不只是“把数组送进模型”；
    \item 学会从波长采样、归一化和缺陷像元开始做最基本预处理；
    \item 用 PCA 先看数据里的主结构；
    \item 用一个中等复杂度分类器完成基线任务；
    \item 回到错分样本，检查问题出在数据、特征还是类别定义本身。
\end{enumerate}

\section{问题背景：为什么光谱是一个理想案例}

在天文和物理里，光谱几乎是最典型的高维数据形式之一。每个样本往往不是两三个特征，而是一整段通量序列；而真正区分类别的，常常是吸收线、发射线、连续谱斜率和信噪比之间的组合。

这让光谱分类成为一个非常好的教学案例，因为它天然包含以下几个层次：

\begin{itemize}
    \item \textbf{预处理问题}：不同样本是否在同一波长网格上，是否需要归一化；
    \item \textbf{特征问题}：直接用全部像元，还是先做压缩表示；
    \item \textbf{分类问题}：使用线性模型、最近邻、随机森林还是更复杂的网络；
    \item \textbf{诊断问题}：模型错分的到底是什么样的样本。
\end{itemize}

因此，这一章不是在“挑一个光谱模型”，而是在训练一种完整的分析顺序：先整理输入，再看结构，再做分类，最后回到个别样本和物理解释。

\section{数据：从一维光谱矩阵开始}

配套 notebook 使用教学数据文件 \texttt{spectral\_classification\_case\_demo.csv}。

它是一个小型 SDSS/LAMOST 风格教学数据集，包含三类对象：

\begin{itemize}
    \item Balmer 吸收主导的恒星；
    \item 金属线更强的冷恒星；
    \item 具有明显 H$\beta$ 和 H$\alpha$ 发射的发射源。
\end{itemize}

数据组织形式是“每个对象一行、每个波长点一列”，再加上 \texttt{snr} 和 \texttt{quality\_flag} 这样的辅助信息。对这个教学集来说，我们只保留 7 个波长采样点，但已经足够演示：

\begin{itemize}
    \item 连续谱斜率；
    \item Balmer 吸收强弱；
    \item H$\beta$/H$\alpha$ 发射增强；
    \item 低信噪比和坏像元如何影响分类。
\end{itemize}

真实项目里，你通常还会同时碰到：

\begin{itemize}
    \item 波长覆盖略有不同的样本；
    \item 某些像元的坏值或缺测；
    \item 不同信噪比；
    \item 来自不同观测夜、不同仪器或不同管线版本的系统差异。
\end{itemize}

所以，光谱分类的第一课通常不是选模型，而是先把输入数据整理成可比状态。

\section{第一步：先做最基本的光谱预处理}

教学案例里，最常见的最小预处理包括：

\begin{itemize}
    \item 对齐或截取到一致的波长区间；
    \item 对坏值做屏蔽或插值；
    \item 用中位数、连续谱或某个参考波段做归一化；
    \item 视情况保留少量物理上清楚的线指标。
\end{itemize}

\begin{examplebox}[一个最小预处理思路]
\begin{lstlisting}[style=pythonstyle]
flux = clean_bad_pixels(raw_flux)
flux = normalize_by_median(flux)
feature_vector = flux
\end{lstlisting}
\end{examplebox}

这里的关键不是“预处理越多越好”，而是要清楚每一步在解决什么问题。归一化是在减少亮度尺度差异，坏值处理是在避免极少数异常像元主导模型，而统一波长网格是在保证样本之间真正可比。

这个案例里就故意放入了一个 \texttt{flux\_4861} 为坏值的样本，用来演示：错误修补并不是无关紧要的小修饰，它会直接影响后续分类位置。

\section{为什么归一化在光谱分类里尤其重要}

如果不做归一化，模型很容易把“整体亮度大小”误当成主要区分信息，而忽略真正关键的谱线形状和相对强弱。对教学光谱来说，这尤其危险，因为我们真正想学的是：

\begin{itemize}
    \item 吸收线有多深；
    \item 发射线有多强；
    \item 连续谱更蓝还是更红；
    \item 不同特征之间是否组合成了某种物理类型。
\end{itemize}

换句话说，归一化不是审美选择，而是在告诉模型：现在你应该更多关注形状，而不是总亮度。

\section{第二步：先用 PCA 看结构，再谈分类}

对高维光谱来说，PCA 往往是一个非常好的第一站。原因很朴素：

\begin{itemize}
    \item 它能告诉你主要变化到底来自哪里；
    \item 它能帮助你判断光谱是否天然形成几类簇；
    \item 它能作为降维手段，为后续分类器提供更稳的低维输入。
\end{itemize}

在很多教学数据里，前几个主成分就已经能分开连续谱温度差异、发射线强弱和某些异常样本。哪怕最后不用 PCA 作为正式特征，这一步也几乎总是值得做。

\begin{examplebox}[PCA 之后最值得先看什么]
\begin{lstlisting}[style=pythonstyle]
pc1, pc2 = projected_spectra[:, 0], projected_spectra[:, 1]
plot(pc1, pc2, color=label)
\end{lstlisting}
\end{examplebox}

\section{PCA 在这个案例里意味着什么}

在本案例中，PCA 不只是为了“降到二维好画图”，而是在帮助我们回答：

\begin{itemize}
    \item 主要变化更像是连续谱斜率差异，还是谱线增强差异；
    \item 正常类群是否已经在前两维里自然分离；
    \item 那些靠近边界或独特的样本是否会在 PCA 平面里突出出来。
\end{itemize}

因此，PCA 在这里扮演的是“结构审查工具”，而不仅是压缩工具。

\section{第三步：建立一个中等复杂度分类 baseline}

这一章最推荐的 baseline 不是一开始就上卷积网络，而是先在 PCA 压缩后的特征上建立一个中等复杂度分类器。配套 notebook 里，我们用的是 \textbf{PCA 前两维上的最近中心分类器}。它足够简单，但已经能回答一个很好的问题：

\begin{quote}
如果主结构真的已经被 PCA 提炼出来，那么只靠低维几何位置，能不能先把大部分对象分开？
\end{quote}

更一般地说，你也可以把这一步替换成：

\begin{itemize}
    \item 线性或核 SVM；
    \item 随机森林；
    \item 最近邻或最近中心模型。
\end{itemize}

选择这类模型有两个教学上的好处：

\begin{itemize}
    \item 训练成本低，便于学生反复实验；
    \item 结果更容易回到 PCA 空间和原始光谱上解释。
\end{itemize}

我们这章强调的是“先把结构看清、把 baseline 做稳”，而不是尽快跳到最复杂的深度学习模型。

\section{第四步：分类报告不够，必须回看错分样本}

光谱分类项目里，最有价值的部分通常不在整体准确率，而在错分样本。你至少应该逐个检查：

\begin{itemize}
    \item 错分样本是否本来就在类别边界附近；
    \item 是否存在低信噪比、坏像元或归一化失败；
    \item 某个类别是不是在训练集中数量太少；
    \item 真实光谱是否同时具有两类对象的特征。
\end{itemize}

在当前教学数据里，\texttt{E005} 就是一个典型例子：它既有坏像元修补，又只有较低信噪比，结果在 PCA 空间里落到了发射源和 Balmer 星之间。这样的样本特别适合做错误案例复查。

这就是为什么本案例的最终产出不只是分类报告，而是\textbf{分类报告加错误案例复查}。在科学场景里，后者往往更重要。

\section{为什么光谱分类特别能训练项目意识}

这一章和前面 Part III 最大的不同，是它迫使你把模型结果放回数据生成过程里理解。例如：

\begin{itemize}
    \item 如果 PCA 的第一主成分几乎只在追踪信噪比，那么分类准确率再高也要警惕；
    \item 如果模型主要依赖少数几个波长点，这些点可能对应物理谱线，也可能只是系统误差；
    \item 如果错分几乎都集中在某个子类，问题可能不在模型，而在标签定义太粗。
\end{itemize}

这就是“项目意识”的核心：模型表现只是结果，数据条件和物理解释才是上下文。

\section{什么时候才应该引入 CNN}

如果基线模型已经建立好、预处理也比较稳定，那么一维卷积网络确实是很自然的下一步。它的优势通常在于：

\begin{itemize}
    \item 能直接从局部线形结构中学习特征；
    \item 不必完全依赖手工压缩后的 PCA 表示；
    \item 在样本足够多时，可能学到更强的模式组合。
\end{itemize}

但这里仍然要保持克制：\textbf{CNN 不是默认答案。} 如果你的样本很小、标签噪声较高、预处理本身还没理顺，那么先把 PCA + 传统分类器这条线跑顺，通常更有教学价值，也更接近真实科研的工作顺序。

\section{这个案例的边界在哪里}

尽管这一章已经很像一个完整的小型项目，但它仍然是教学版，而不是正式的大巡天光谱管线。它省略了很多真实复杂性，例如：

\begin{itemize}
    \item 更密集的波长采样和更高维的特征空间；
    \item 天地线残留、通量定标误差、拼接问题等系统效应；
    \item 更复杂的类别体系与混合天体；
    \item 数据集规模扩大后对模型复杂度和计算成本的要求。
\end{itemize}

因此，我们不应把这个案例理解成“光谱分类已经被解决”，而应把它理解成“怎样从原始光谱矩阵走到一个可解释分类结论”的最小练习场。

\section{一个最小光谱分类案例报告模板}

光谱分类报告最容易写偏成“某分类器准确率是多少”。更好的写法，是把预处理、结构观察、分类结果和错分复查连起来。一个最小报告可以包含：

\begin{enumerate}
    \item \textbf{光谱输入契约}：说明波长网格、通量列、信噪比字段和质量标记；
    \item \textbf{预处理记录}：说明坏像元如何处理、归一化使用什么参考、是否剔除低质量样本；
    \item \textbf{PCA 结构}：解释前两维主要对应连续谱斜率、吸收/发射线强度，还是其他变化；
    \item \textbf{分类 baseline}：报告模型、特征空间、混淆矩阵和主要类别指标；
    \item \textbf{错分光谱复查}：挑出最典型错分样本，回到原始光谱、SNR 和坏像元记录；
    \item \textbf{升级条件}：说明何时才有必要引入一维 CNN 或更复杂表示学习。
\end{enumerate}

\begin{examplebox}[光谱分类案例报告骨架]
\begin{lstlisting}[style=pythonstyle]
spectral_report = {
    "preprocessing": ["bad-pixel repair", "median normalization"],
    "pca_interpretation": pca_component_notes,
    "classifier": "nearest centroid on first two PCs",
    "confusion_matrix": confusion_matrix,
    "misclassified_spectra": review_ids,
    "cnn_upgrade_reason": "only if baseline fails on stable local line patterns",
}
\end{lstlisting}
\end{examplebox}

这样的报告能防止一个常见误读：把光谱分类当成普通表格分类。光谱的每一列都来自波长、仪器和物理线形，分类结论必须能回到这些结构中解释。

\begin{warnbox}[光谱分类中的常见误区]
\begin{itemize}
    \item 没有统一归一化方式，就直接比较不同光谱；
    \item 一上来就训练复杂网络，却没有先做 PCA 可视化；
    \item 只看整体准确率，不检查错分对象的光谱形状；
    \item 把仪器系统差异误当作天体物理差异；
    \item 样本量很小时仍然过度追求模型复杂度。
\end{itemize}
\end{warnbox}

\section{AI 助手如何帮助你}

在这个案例里，AI 助手适合帮助你：

\begin{itemize}
    \item 检查预处理、PCA 和分类代码是否自洽；
    \item 帮你把分类报告和错误案例表整理得更清楚；
    \item 梳理“错分原因可能来自哪里”的分析框架；
    \item 协助你把 PCA 平面结构和原始光谱形状放到同一套解释里。
\end{itemize}

但某条光谱特征究竟对应真实物理结构还是数据管线问题、哪个错分案例值得进一步跟进、以及是否应该升级到 CNN，依然需要你自己判断。

\section{练习}

\begin{exercisebox}[基础练习]
在 notebook 中关闭坏像元修补步骤，直接用原始 \texttt{flux\_4861} 进入归一化和 PCA。比较 PCA 平面、分类结果和错分样本是否改变，并说明一个坏像元为什么可能影响整条光谱的位置。
\end{exercisebox}

\begin{exercisebox}[进阶练习]
把归一化方法从中位数归一化改成用某个参考波段归一化，重新比较 PCA 前两维和最近中心分类结果。请解释：不同归一化方式是在强调整体谱形、局部线强，还是某个波段的相对比例。
\end{exercisebox}

\begin{exercisebox}[开放练习]
设计一个把本章 workflow 扩展到真实巡天光谱的最小方案。至少包含：统一波长网格、屏蔽坏像元、信噪比阈值、训练/测试划分、错分光谱复查，以及何时才值得引入一维 CNN。
\end{exercisebox}

\section{本章小结}

光谱分类案例最重要的价值，不在于把某个分类器调到多高分，而在于训练一种完整工作方式：先整理输入、再看结构、再做 baseline、最后回到错分样本和物理解释。只有这样，分类结果才真正有科学意义。本章真正训练的是项目意识，而不是模型崇拜。
